\chapter{The shape of every raylib program}

\section{Immediate mode, and why nothing stays on screen}

Here is the idea that everything else in raylib hangs from. Most graphical
toolkits you have met are \emph{retained mode}: you create a button object, you
add it to a window, and it stays there. The toolkit remembers it and redraws it
for you.

raylib is \emph{immediate mode}. Nothing is remembered. There is no circle
object. \lstinline|DrawCircle()| does not create a circle --- it paints circle
shaped pixels, right now, and then forgets all about it. If you want the circle
to still be there in a sixtieth of a second, you call \lstinline|DrawCircle()|
again.

\keyidea{Your program repaints the entire picture, from a blank canvas,
sixty times a second. Everything you can see on screen is something you drew
this frame. Everything you stop drawing disappears instantly.}

This sounds wasteful and it is not: GPUs are extremely good at exactly this. And
it buys you something precious for a beginner --- there is no hidden state. If
something is on screen, a line in your code put it there this frame. If something
is missing, a line in your code did not run. Debugging graphics becomes reading
your own draw calls in order.

\section{The skeleton}

Every raylib program, from the ten line example to the full game, has this shape:

\begin{lstlisting}
#include "raylib.h"

int main(void)
{
    InitWindow(800, 450, "My game");   // 1. once: open the window
    SetTargetFPS(60);

    while (!WindowShouldClose())       // 2. the frame loop
    {
        // UPDATE: change variables, draw nothing

        BeginDrawing();                // 3. start painting
            ClearBackground(RAYWHITE);
            DrawText("Hello", 20, 20, 20, DARKGRAY);
        EndDrawing();                  // 4. show it
    }

    CloseWindow();                     // 5. once: clean up
    return 0;
}
\end{lstlisting}

That is lesson 1 of the 2D course, near enough. Let us take the five parts apart.

\begin{center}
\begin{tikzpicture}[
  node distance=8mm,
  every node/.style={font=\sffamily\small},
  bx/.style={draw=codeframe,fill=rlight,rounded corners=2pt,minimum width=6.2cm,minimum height=8mm,align=center},
  lbx/.style={bx,fill=white,draw=rblue!40},
  arrow/.style={-{Stealth[length=5pt]},rgrey}
]
\node[bx] (init) {\texttt{InitWindow()} \quad once};
\node[lbx,below=of init] (check) {\texttt{WindowShouldClose()} ?};
\node[lbx,below=of check] (update) {update your variables};
\node[lbx,below=of update] (begin) {\texttt{BeginDrawing()}};
\node[lbx,below=of begin] (clear) {\texttt{ClearBackground()}};
\node[lbx,below=of clear] (draw) {all your \texttt{Draw...()} calls};
\node[lbx,below=of draw] (end) {\texttt{EndDrawing()} \ \ swap buffers};
\node[bx,below=14mm of end] (close) {\texttt{CloseWindow()} \quad once};

\draw[arrow] (init) -- (check);
\draw[arrow] (check) -- node[right,font=\sffamily\tiny,text=rgrey] {no} (update);
\draw[arrow] (update) -- (begin);
\draw[arrow] (begin) -- (clear);
\draw[arrow] (clear) -- (draw);
\draw[arrow] (draw) -- (end);
\draw[arrow] (end.east) -- ++(1.5,0) |- node[right,pos=0.25,font=\sffamily\tiny,text=rgrey] {next frame} (check.east);
\draw[arrow] (end) -- node[right,font=\sffamily\tiny,text=rgrey] {yes, quit} (close);
\end{tikzpicture}
\end{center}

\subsection{1. InitWindow}

\lstinline|InitWindow(width, height, title)| creates the window and, more
importantly, the OpenGL context behind it. Until this line runs, no other raylib
function works --- not loading a texture, not measuring text, nothing. Calling
\lstinline|LoadTexture()| before \lstinline|InitWindow()| is a classic
first-day crash.

Flags are the exception: \lstinline|SetConfigFlags()| must be called
\emph{before} \lstinline|InitWindow()|, because it configures how the window
gets created. The two you will want:

\begin{lstlisting}
SetConfigFlags(FLAG_VSYNC_HINT | FLAG_MSAA_4X_HINT);
InitWindow(800, 450, "My game");
\end{lstlisting}

\lstinline|FLAG_VSYNC_HINT| syncs your drawing to the monitor's refresh so the
image does not tear. \lstinline|FLAG_MSAA_4X_HINT| turns on antialiasing so
diagonal edges stop looking like staircases.

\subsection{2. WindowShouldClose}

\lstinline|while (!WindowShouldClose())| is your game loop. The function returns
true when the user presses Escape or clicks the window's close button, so the
loop reads as ``keep going until they want out''.

Each pass through this loop is one \emph{frame}. At 60 FPS the body of that loop
runs 60 times a second, which gives you about 16 milliseconds to do everything.

\subsection{3 and 4. BeginDrawing and EndDrawing}

This pair is where the confusion usually is, so let us be precise.

\lstinline|BeginDrawing()| says ``I am about to paint''. It prepares the canvas
for this frame. Every \lstinline|Draw...()| call after it adds geometry to a
batch that raylib is accumulating.

\lstinline|EndDrawing()| does three jobs at once:

\begin{enumerate}[leftmargin=*,itemsep=2pt]
  \item it flushes the batch --- everything you queued actually gets sent to the
        GPU and drawn;
  \item it swaps the buffers. You have been painting on a hidden canvas the whole
        time; this is the moment it becomes the visible one. That is
        \emph{double buffering}, and it is why you never see a half-drawn frame;
  \item it polls input events and waits, if needed, to hit your target frame rate.
\end{enumerate}

\gotcha{Nothing appears on screen until \lstinline|EndDrawing()| runs. If you
put a \lstinline|sleep()| or a long loop between \lstinline|BeginDrawing()| and
\lstinline|EndDrawing()| hoping to see an animation, you will see nothing at all
but a frozen window. Animation happens \emph{across} frames, never inside one.}

\subsection{ClearBackground}

The window still holds last frame's pixels. \lstinline|ClearBackground(RAYWHITE)|
wipes them. Skip it and you get smears --- which, incidentally, is a fun thing to
try once on purpose: comment it out in lesson 1 and move something around.

\section{Update and draw: keep them apart}

Notice how the skeleton has a comment saying ``UPDATE'' before
\lstinline|BeginDrawing()|. That separation is not enforced by raylib, it is a
discipline, and it is worth adopting from day one:

\begin{itemize}[leftmargin=*,itemsep=2pt]
  \item \textbf{Update} --- read input, move things, run the AI, check
        collisions. Change variables. Draw nothing.
  \item \textbf{Draw} --- paint the current state of those variables. Change nothing.
\end{itemize}

When the two get tangled you end up with bugs where an object reacts to a
collision that was resolved half a frame ago, and they are miserable to find.
Keeping them apart costs nothing and every lesson in this course does it.

\section{Frame rate, and the two ways to control it}

\lstinline|SetTargetFPS(60)| tells raylib to throttle the loop to at most 60
iterations per second. Without it your loop runs as fast as the machine allows,
burning 100\% of a CPU core to draw frames nobody will see.

\lstinline|FLAG_VSYNC_HINT| does a similar job at a different level: it makes the
buffer swap wait for the monitor. In practice, set both; they cooperate.

Two useful readings, both free:

\begin{lstlisting}
int   fps = GetFPS();          // frames per second, as an integer
float dt  = GetFrameTime();    // seconds the LAST frame took: ~0.0166 at 60fps
double t  = GetTime();         // seconds since InitWindow()
\end{lstlisting}

\lstinline|GetFrameTime()| is the important one and it gets its own chapter
(chapter 5). It is the difference between a game that behaves the same on every
computer and one that does not.

\tryit{Lesson 1 prints the frame counter, the elapsed time and the frame
duration, live. Watch the frame counter climb and realise that every one of those
numbers was repainted from scratch.}{./course2d/build.sh 1}
